\metatitle{Chromatic symmetric functions in the elementary basis}
\metadescription{A guide to elementary-basis expansions of chromatic symmetric
and chromatic quasisymmetric functions, including the Stanley--Stembridge and
Shareshian--Wachs conjectures, q-refined results,
recursive methods, and known negative examples.}
\metakeywords{chromatic symmetric functions,e-positivity,Stanley--Stembridge conjecture,
Shareshian--Wachs conjecture,unit interval graphs}

\section[chromaticEexpansion]{Chromatic symmetric functions in the elementary basis}

Let $G$ be a graph and let $\chrom_G(\xvec)$ denote its
\hyperref[chromaticQuasisymmetricDefinition]{chromatic symmetric function}.
An \defin{$\elementaryE$-expansion} is an expression
\label{chromaticEExpansionDefinition}
\[
\chrom_G(\xvec) = \sum_{\lambda} c_\lambda \elementaryE_\lambda(\xvec).
\]
We say that $\chrom_G(\xvec)$ is
\defin{$\elementaryE$-positive}\label{chromaticEPositive} if all
$c_\lambda$ are non-negative integers. If $G$ is naturally labeled and
$\chrom_G(\xvec;q)$ is symmetric, then the stronger $q$-refined question asks
whether
\[
\chrom_G(\xvec;q) = \sum_{\lambda} c_\lambda(q) \elementaryE_\lambda(\xvec)
\qquad\text{with}\qquad c_\lambda(q) \in \setN[q].
\]

The $q$-refined question is only available in families where
$\chrom_G(\xvec;q)$ is symmetric, most importantly for
\hyperref[unitIntervalGraph]{unit interval graphs} or
\hyperref[circularUnitArcDigraph]{circular unit arc digraphs}.
For general graphs, the default question is the specialization $q=1$.


\subsection[chromaticEExpansionConjectures]{The two main conjectures}

The \hyperref[stanleyStembridgeConjecture]{Stanley--Stembridge conjecture}
was formulated by \name{Richard Stanley} and \name{John Stembridge}
\cite{StanleyStembridge1993}, and then in the language of chromatic symmetric
functions by Stanley \cite{Stanley1995}.

\begin{theorem}[Hikita, \cite{Hikita2024x}]
Let $P$ be a $(3+1)$-free poset. Then the chromatic symmetric function of
the incomparability graph of $P$ is $\elementaryE$-positive.
\end{theorem}

By work of \name{Mathieu Guay-Paquet} \cite{GuayPaquet2013}, it is enough to
prove the result for unit interval graphs. Hikita's proof gives a formula for
the elementary coefficients of the Shareshian--Wachs
$q$-chromatic refinement, and this formula specializes positively at $q=1$.
Thus the theorem settles the Stanley--Stembridge conjecture, but it does not
settle the following refinement.

Hikita later introduced \defin{$(q,t)$-chromatic symmetric functions}
\label{qtChromaticSymmetricFunction} for
unit interval graphs using level-one polynomial representations of affine
Hecke algebras of type $A$ \cite{Hikita2025x}.  These functions refine the
chromatic quasisymmetric functions in the sense that Hikita recovers the
Shareshian--Wachs function from the $(q,t)$-theory
\cite[Thm.~6.1]{Hikita2025x}.  The specialization at $q=\infty$ gives the
probabilistic interpretation of the elementary coefficients used in Hikita's
proof of the Stanley--Stembridge conjecture.
\name{JiSun Huh}, \name{Byung-Hak Hwang}, \name{Donghyun Kim},
\name{Jang Soo Kim}, and \name{Jaeseong Oh} refine Hikita's theorem using
the $g_{\mathfrak{m},k}(\xvec;q)$-functions of
\name{Alex Abreu} and \name{Antonio Nigro} \cite{HuhHwangKimKimOh2025x}.
They prove the $\elementaryE$-positivity of
$g_{\mathfrak{m},k}(\xvec;1)$, give a Schur expansion for a related
weighted sum in terms of $P$-tableaux with $1$ in the upper-left corner, and
introduce a restricted modular law which determines such functions from
their values on disjoint unions of paths.

\begin{conjecture}[Shareshian--Wachs, \cite[Conj.~4.9]{ShareshianWachs2012}]
\label{chromaticEShareshianWachsConjecture}
If $G$ is a unit interval graph with its natural labeling, then
\[
\chrom_G(\xvec;q) = \sum_{\lambda} c_\lambda(q)\elementaryE_\lambda(\xvec)
\qquad\text{has}\qquad c_\lambda(q) \in \setN[q].
\]
\end{conjecture}

A desirable combinatorial proof would give a statistic $\mu$ on acyclic
orientations of $G$ such that
\[
\chrom_G(\xvec;q)
= \sum_{\theta \in AO(G)} q^{\asc(\theta)}
\elementaryE_{\mu(\theta)}(\xvec).
\]
This is the form in which one can see both the $q$-grading and the
elementary-basis positivity at the same time.
\name[Gergely Berczi]{Gergely B{\'e}rczi} and
\name[Jonas Kluver]{Jonas Kl{\"u}ver} use reinforcement learning to propose
a universal counting formula for coefficients of chromatic symmetric
functions of unit interval graphs \cite{BercziKluver2024x}.  Their
conjectural formula counts disjoint tuples of Eschers satisfying
graph-independent concatenation conditions.
The conjecture extends to proper circular unit arc digraphs; see
\cite{Ellzey2017,AlexanderssonPanova2018}.


\begin{conjecture}[See \cite[Conj.~5.1]{ShareshianWachs2016}]
\label{chromaticEUnimodalityConjecture}
Write
\[
\chrom_G(\xvec;q) = \sum_{j=0}^m q^j a_j(\xvec),
\]
where $m$ is the number of edges of a unit interval graph $G$.
Then $a_{j+1}(\xvec)-a_j(\xvec)$ is $\elementaryE$-positive for
all $0 \leq j \lt (m-1)/2$.
\end{conjecture}

This unimodality conjecture also extends to the
\hyperref[circularUnitArcDigraph]{circular unit arc digraph} setting.

\subsection[chromaticEExpansionMethods]{Kinds of elementary expansions}

It is useful to distinguish the following kinds of results.

\begin{itemize}
\item
A \emph{combinatorial} formula gives the coefficient of
$\elementaryE_\lambda$ as the cardinality or weight-enumerator of a set of
objects.  This is the strongest kind of evidence for an eventual
Shareshian--Wachs proof.

\item
A \emph{sign-reversing involution} starts from a signed
$\elementaryE$-expansion and cancels all negative terms. Several proofs for
special unit interval families have this form.

\item
A \emph{recursive} or generating-function proof shows that the chromatic
symmetric functions satisfy an $\elementaryE$-positive recurrence with
$\elementaryE$-positive initial data. This proves positivity, but does not
always identify a natural coefficient statistic.

\item
A \emph{noncommutative} proof establishes the stronger
$(\elementaryE)$-positivity of the chromatic symmetric function in
\hyperref[schurInNonCom]{noncommuting variables}, as in
\cite{GebhardSagan2001}.  This implies ordinary
$\elementaryE$-positivity.

\item
A \emph{negative} result usually means that one graph in the stated family has
a negative elementary coefficient. It does not mean that every graph in that
family is non-positive.
\end{itemize}

\subsection[chromaticEExpansionChronology]{Chronological guide}

The following list is a compact guide to the main families and methods.

\begin{itemize}
\item
\emph{1993--1995.}
Stanley and Stembridge introduce the conjectural positivity problem
\cite{StanleyStembridge1993}. Stanley proves $\elementaryE$-positivity for
paths and cycles, and obtains positive formulas for complete graphs and
co-triangle-free graphs \cite{Stanley1995}.

\item
\emph{2001.}
\name{David Gebhard} and \name{Bruce Sagan} introduce a chromatic symmetric
function in noncommuting variables and prove noncommutative
$(\elementaryE)$-positivity for $K_\alpha$-chains and for diamond and path
chains \cite{GebhardSagan2001}.

\item
\emph{2012--2016.}
\name{John Shareshian} and \name{Michelle Wachs} introduce the
$q$-refinement and the Shareshian--Wachs conjecture
\cite{ShareshianWachs2012,ShareshianWachs2016}.  The $q$-refinement is
symmetric for unit interval graphs and specializes to the ordinary chromatic
symmetric function at $q=1$.

\item
\emph{2017--2022.}
Several special cases for unit interval graphs are proved: triangular ladders
\cite{Dahlberg2018x}, lollipops and lariats \cite{DahlbergWilligenburg2018},
the Abelian and bounce-number-two cases
\cite{HaradaPrecup2019,ChoHuh2019,NadeauTewari2022x,LeeSoh2022x}, and
melting lollipop families \cite{HuhNamYoo2020}.

\item
\emph{2018--2021.}
For graph classes outside the unit interval setting, positive and negative
families are sorted in work on $H$-free graphs
\cite{Tsujie2018,HamelHoangTuero2019,FoleyHoangMerkel2019}, on
claw-contractible-free graphs \cite{DahlbergFoleyWilligenburg2020}, and on
$(claw,2K_2)$-free graphs \cite{LiYang2021}.

\item
\emph{2021--2023.}
Certain cycle-chord and tadpole-type families are treated using
noncommutative and recursive methods \cite{WangWang2022}.
\name{Foster Tom} gives a signed
$\elementaryE$-expansion which yields new positive graph families
\cite{Tom2023x}.  Coefficient-level positivity for partitions with at most two
parts is proved in \cite{AbreuNigro2023x,RokSzenes2023x}.

\item
\emph{2024--2026.}
Recent results include the $2+1+1$-avoiding
\hyperref[unitIntervalOrder]{unit interval orders}
\cite{McDonoughPylyavskyyWang2024x}, twinned paths and cycles
\cite{BanaianCelanoChangLeeColmenarejoGoffKimbleKimpelLentferLiangSundaram2025},
all cycle-chord graphs \cite{Wang2025}, clocks \cite{ChenHeWang2026},
conjoined graphs \cite{QiTangWang2025}, gluing graphs at a single vertex
\cite{TomVailaya2025x}, adjacent cycle-chains \cite{TomVailayaAdjacent2026},
twinned lollipops and kayak paddle graphs \cite{TangWang2024x}, and
refinements of Hikita's proof
\cite{HuhHwangKimKimOh2025x,GriffinMellitRomeroWeiglWen2025x}.
Hook-shape immanant characters can be written as non-negative sums of
Stanley--Stembridge characters by \name{Nathan R. T. Lesnevich}
\cite{Lesnevich2024}; this proves a hook-shape case of a conjecture of
\name{Richard Stanley} and \name{John Stembridge} on immanant characters.
\end{itemize}


\subsection[chromaticEExpansionQRefined]{Unit interval and $q$-refined results}

When the graph is a unit interval graph, one should record whether the result
proves the full $q$-refined statement or only the specialization $q=1$.
The following families are among the main cases where the refinement is known
or where the proof is naturally stated in the chromatic quasisymmetric setting.

\begin{itemize}
\item
Complete graphs, paths, cycles, and the directed-cycle analogues have explicit
formulas in the Shareshian--Wachs/Ellzey setting
\cite{ShareshianWachs2012,Ellzey2017,AlexanderssonPanova2018}.

\item
If both a unit interval graph and its complement are unit interval graphs, then
$\chrom_G(\xvec)$ is $\elementaryE$-positive
\cite{FoleyHoangMerkel2019}.

\item
For Abelian area sequences and related bounce-number-two families, the
literature contains recursive, cohomological, and sign-reversing-involution
proofs \cite{HaradaPrecup2019,ChoHuh2019,NadeauTewari2022x,LeeSoh2022x}.

\item
For area sequences with bounce number three, some coefficients are treated in
\cite{ChoHong2019}, and this is extended in \cite{Wang2022x}.

\item
All elementary coefficients indexed by partitions with at most two parts are
non-negative, by \cite{AbreuNigro2023x} using the cohomology of
\hyperref[hessenbergVarieties]{Hessenberg varieties}, and independently by
\cite{RokSzenes2023x}.

\item
The $2+1+1$-avoiding \hyperref[unitIntervalOrder]{unit interval orders},
equivalently the area sequences with $i-2 \leq a_i$, are treated in
\cite{McDonoughPylyavskyyWang2024x} using strand diagrams and
\hyperref[representationTheory]{representation theory}.
\end{itemize}



\subsection[chromaticEExpansionOrdinaryFamilies]{Further positive families}

The following examples are useful landmarks for the ordinary $q=1$ problem.

\begin{itemize}
\item
Noncommutative $(\elementaryE)$-positive families include $K_\alpha$-chains,
diamond and path chains, and several later families proved by adapting the
Gebhard--Sagan framework \cite{GebhardSagan2001,WangWang2022,WangZhou2024x}.

\item
Lollipops, lariats, triangular ladders, and generalized pyramid graphs give
some of the standard small graph families where direct signed expansions can
be made positive \cite{Dahlberg2018x,DahlbergWilligenburg2018,LiYang2021}.

\item
Twinning preserves $\elementaryE$-positivity for paths and cycles, with both
positive generating functions and positive recurrences
\cite{BanaianCelanoChangLeeColmenarejoGoffKimbleKimpelLentferLiangSundaram2025}.
Twinning does not preserve $\elementaryE$-positivity for arbitrary graphs
\cite{LiLiWangYang2021}.

\item
Cycle-chord graphs are $\elementaryE$-positive
\cite{WangWang2022,Wang2025}.  The proof of the general case uses the
composition method of \cite{WangZhou2024x}.

\item
Graphs obtained by gluing at a single vertex give a flexible source of
new examples.  In particular, gluing sequences of unit interval graphs and
cycles gives $\elementaryE$-positive graphs \cite{TomVailaya2025x}.

\item
Adjacent cycle-chains are $\elementaryE$-positive, and the same methods
extend to graphs formed by connecting a sequence of cycles and cliques
\cite{TomVailayaAdjacent2026}.

\item
Positive $\elementaryE$-expansions are known for KPKP graphs, twinned
lollipops, and kayak paddle graphs \cite{TangWang2024x}.  This refines the
earlier $\elementaryE$-positivity of kayak paddle graphs from
\cite{AliniaeifardWangWilligenburg2024}.

\item
Several tree families are now classified. For example, \cite{WangWang2023}
classifies positivity for all broom graphs and most double broom graphs,
and \cite{TangWangWang2024x} proves $\elementaryE$-positivity for the
spiders $S(4m+2,2m,1)$.
\end{itemize}


\subsection[chromaticEExpansionNegative]{Negative results and obstructions}

Negative results are important: they indicate which graph-theoretic hypotheses
are doing real work.

\begin{itemize}
\item
The star $K_{1,3}$ is already not $\elementaryE$-positive.  More generally,
large-degree cut vertices and several families of trees obstruct
$\elementaryE$-positivity \cite{DahlbergSheWilligenburg2020}.

\item
The saltire, augmented saltire, and triangular tower families give negative
examples in the study of claw-contractible-free graphs
\cite{DahlbergFoleyWilligenburg2020}.

\item
For $H$-free graph classes, some families are positive and some contain
counterexamples.  See \cite{HamelHoangTuero2019,FoleyHoangMerkel2019} for a
systematic treatment.

\item
Some spider and related families are not $\elementaryE$-positive; see
\cite{FoleyKazdanKrollAlbergaMelnykTenenbaum2018x} and the later work on
trees and connected partitions \cite{Tom2026}.

\item
Tom proves that trees with a vertex of degree at least $5$, trees with a
degree-$4$ vertex not adjacent to a leaf, and four-legged spiders are not
$\elementaryE$-positive \cite{Tom2026}.  A quantitative connected-partition
approach to further tree obstructions is developed in \cite{EthanLi2025x}.

\item
Twinning an $\elementaryE$-positive graph at a vertex can destroy even Schur
positivity, and hence can destroy $\elementaryE$-positivity
\cite{LiLiWangYang2021}.
\end{itemize}

\begin{example}[The smallest star obstruction]\label{chromaticEStarObstruction}
For the star graph $K_{1,3}$,
\[
\chrom_{K_{1,3}}(\xvec)
= \elementaryE_{211} + 5\elementaryE_{31}
-2\elementaryE_{22} + 4\elementaryE_4.
\]
The coefficient of $\elementaryE_{22}$ is negative, so this graph is not
$\elementaryE$-positive.
\end{example}


\subsection[chromaticEExpansionCoefficientTests]{Coefficient-level tests}

Several papers study individual elementary coefficients rather than a whole
family of graphs.  For example, \cite{ChoHongLee2020x} gives explicit
conditions for when a coefficient in the $\elementaryE$-expansion is positive,
and \cite{CrewZhang2022x} studies elementary-basis coefficients in relation to
\hyperref[acyclicOrientation]{acyclic orientations} and sinks.  The
change-of-basis perspective in
\cite{SaganTom2026} gives another way to turn information about particular
partitions into graph-theoretic positivity statements.

After Hikita's proof, several papers have focused on making the coefficients
more explicit.  The Macdonald expansion of
\cite{GriffinMellitRomeroWeiglWen2025x} rederives Hikita's formula, while
\cite{Siegl2025x} studies upper and lower combinatorial bounds for elementary
coefficients.  Finally, \cite{Kravitz2026x} proves that the partitions which
always have non-negative elementary coefficients for every finite graph are
precisely the hook partitions.
